% Generated draft: requires mathematical and source review.
\documentclass[11pt]{article}
\usepackage[margin=25mm]{geometry}
\usepackage{fontspec}
\setmainfont{Latin Modern Roman}
\usepackage{amsmath,amssymb,mathtools,unicode-math}
\setmathfont{Latin Modern Math}
\usepackage{enumitem,needspace,xurl,hyperref}
\setlength{\parindent}{0pt}
\setlength{\parskip}{5pt}
\setlength{\emergencystretch}{3em}
\newcommand{\N}{\mathbb{N}}
\newcommand{\Z}{\mathbb{Z}}
\newcommand{\Q}{\mathbb{Q}}
\newcommand{\R}{\mathbb{R}}
\newcommand{\C}{\mathbb{C}}
\newcommand{\F}{\mathbb{F}}
\newcommand{\A}{\mathbb{A}}
\newcommand{\PP}{\mathbb{P}}
\newcommand{\E}{\mathbb{E}}
\newcommand{\eps}{\varepsilon}
\newcommand{\dd}{\,\mathrm{d}}
\begin{document}
\section*{UnsolvedMath ID: 11000144}
\textbf{Source problem code:} catalogue entry
\subsection{Definitions and mathematical statement}
The supplied statement of the Ehrenpreis Conjecture says that two closed Riemann surfaces have finite unbranched covers with homeomorphic total spaces arbitrarily close in the Teichm\"uller metric. The usual hyperbolic formulation requires both base genera to be at least $2$, a hypothesis missing from this excerpt.

With that qualification, let $X_1,X_2$ be closed Riemann surfaces of genera $g_1,g_2\geq2$. For every $\varepsilon>0$, the assertion is that there exist connected finite unbranched holomorphic covers $p_i:Y_i\to X_i$ and an orientation-preserving quasiconformal homeomorphism $h:Y_1\to Y_2$ such that
\[
 \tfrac12\log K(h)<\varepsilon,
\]
where $K(h)$ is maximal dilatation. Equivalently, after using an appropriate marking to identify their underlying surfaces, the total-space structures have Teichm\"uller distance less than $\varepsilon$. Covering degrees and common genus may depend on $\varepsilon$. If the degrees are $d_1,d_2$, homeomorphism requires $d_1(g_1-1)=d_2(g_2-1)$; equal covering degrees are not required. The covers need not be conformally identical or be one common holomorphic cover.

The literal unrestricted phrase ``two closed Riemann surfaces'' includes incompatible genus types, for example sphere and positive-genus surfaces, that cannot have homeomorphic finite unbranched covers. Thus the hyperbolic hypothesis must be checked rather than silently omitted. The registry marks the conjecture solved in the literature by Kahn--Markovic; this is supplied status information distinct from the older source's problem label.
\subsection{Short English statement}
For any two closed hyperbolic Riemann surfaces, finite unbranched covers can be made arbitrarily close after choosing a homeomorphism of their total spaces. The missing hyperbolic-genus hypothesis and supplied solution citation require review.
\subsection{Sources}
\begin{itemize}
\item Local registry: \url{data/unsolvedmath/problems.json}, permanent UnsolvedMath ID $11000144$, exact source code \texttt{catalogue entry}.
\item Farb, \emph{Problems on Mapping Class Groups and Related Topics}, Problem 15, PDF page 123: \url{https://www.math.uchicago.edu/~farb/papers/mcgbook.pdf}. Supplied access date: 2026-07-29; extraction: pdf-font-structure; repeated local numbering disambiguated by page among 329 headings.
\item The registry attributes a proof to Kahn--Markovic (2015), giving \emph{J. Amer. Math. Soc.} 28 (2015), 1185--1210. It also mentions 2011 work described as ``The bow lemma,'' ``Immersing almost geodesic surfaces in a closed hyperbolic three manifold,'' and ``Counting essential surfaces.'' These are supplied bibliographic claims; their precise matching to the Ehrenpreis proof is not verified here.
\end{itemize}
\end{document}
